\section{把 AI 变成写作搭档}
\subsection{先连接你自己的模型}
首次进入时，如果你尚未配置 API，网页会弹出引导。你也可以随时前往“设置 → 个人 AI 配置”。选择接口类型，填写服务提供方给出的 API 地址、API Key 和准确的模型名称，然后保存。自定义服务通常使用“OpenAI 兼容接口”。

例如，地址可以是 \url{https://api.example.com/v1}。这是格式示例，并不是一个可用的模型服务。兼容接口也可以粘贴完整的 \texttt{/v1/chat/completions} 地址，保存时会自动整理。模型名称必须使用你的服务实际支持的名称，不能只照抄示例。

密钥由服务器加密保存，普通用户之间互不可读；服务器管理员仍管理服务器上的加密凭据。修改已有配置时，API Key 留空会保留原密钥。使用哪家服务、模型是否支持图片、调用费用如何计算，都取决于你自己填写的模型服务。不要把 API Key 写进文档或对话。

\subsection{选中文本，开始多轮对话}
在编辑器中选中一段正文，会出现“选区对话”浮窗。也可以选中后按快捷键：Mac 使用 Command+K，其他系统使用 Ctrl+K。输入要求后发送，例如：
\begin{itemize}
\item “润色这段，保留原意和 LaTeX 命令。”
\item “让这段更简洁，但不要省略条件。”
\item “先解释这段的逻辑，不要修改正文。”
\end{itemize}
你可以继续追问：“上一版太正式了，换成自然一些的表达。” 当前窗口保留本轮对话，关闭窗口或切换文件后会开始新的对话。

\subsection{先看差异，再决定是否采用}
提出修改时，浮窗会展示删除和新增的内容。点击“接受修改”才会把方案写入选区；点击“拒绝”会保留原文。已经接受的本次修改可以通过“撤销 AI 修改”还原。点击“停止”会取消本轮生成，未完成的内容不会写入文档。

如果你或协作者在生成过程中修改了原选区，系统会暂停接受与撤销，避免覆盖新内容。请重新选择目标段落，开始新的对话。最终仍需你检查事实、公式、引用和论证是否正确。

\subsection{让模型参考项目里的文件和图片}
浮窗中默认勾选“允许读取当前项目文件与图片”。开启时，模型会收到当前文件的编辑快照与项目目录，并可按需调用只读的列目录、读文件工具。状态栏会显示正在做什么，下面会列出本轮已经读取的文件名。

可以尝试：\textbf{“先查看项目目录，阅读其他章节和示例图片，再解释这个项目的写作流程。只回答，不修改正文。”}

这些工具只读当前账号有权限访问的项目，不执行服务器命令。项目内容会发送给你配置的模型服务。取消勾选后，本轮只发送选区、相邻少量上下文和对话。

目前支持 UTF-8 文本、LaTeX、BibTeX、Markdown 和 PNG/JPEG/WebP 图片。图片需要模型支持视觉；PDF、Word、HEIC 和压缩包暂不解析。过长文本会截断，并把截断情况告知模型；工具调用次数也有限制。需要大范围分析时，可以分章节提问。

\subsection{保留自己的判断}
文件被读取，并不意味着模型的回答一定正确。要求引用文献时，让模型使用项目中已有的 BibTeX 键，并核查来源。对于图片，让它说明实际能看到的内容。不要把推测当作事实，也不要把生成的示例数字直接用于论文结论。
